\section{Erd\H{o}s Problem \#311 (Round 2)}

\subsection*{1) ROUND-2 OBJECTIVE}
\textbf{Path C (correction / exact equivalence of formulations).}

The biggest Round-1 gap was that the write-up used the original Erd\H{o}s--Graham formulation with the extra ``no subsubset sums to $1$'' condition, whereas the modern webpage uses a simpler formulation without that condition. The most valuable Round-2 move is therefore to prove that the two formulations are in fact \emph{exactly equal for every $N$}, not merely asymptotically equivalent. Once that is done, the problem can be attacked in the cleaner simplified form, and exact computations become easier to interpret.

\subsection*{2) Round-1 FOUNDATION USED}
I use the following Round-1 results.

\begin{itemize}
\item The denominator lower bound
\[
\delta(N)\ge \frac{1}{[1,\ldots,N]}.
\]
\item Exact values for the original Erd\H{o}s--Graham quantity up to $N=25$.
\item The meet-in-the-middle scaling principle: after multiplying by $L_N:=[1,\ldots,N]$, reciprocal-sum questions become exact integer subset-sum questions with weights $L_N/n$.
\end{itemize}

\subsection*{3) NEW INSIGHT / TOOL (ROUND-2)}
The new idea is a very simple threshold principle:

\medskip
\centerline{\emph{Any set $A\subseteq\{1,\ldots,N\}$ with a strict subsubset summing to $1$ has nonzero distance from $1$ at least $1/N$.}}
\medskip

At the same time, for every $N\ge 3$ there is an explicit subset whose distance from $1$ is \emph{strictly smaller} than $1/N$. Therefore the actual minimizer can never contain such a strict $1$-summing subsubset, so the two formulations are automatically identical.

This removes the awkward condition completely, with no brute-force small-$N$ check.

\subsection*{4) ATTACK PLAN (ROUND-2)}
The exact Round-1 gap was the unproved equivalence between the original and simplified definitions. I therefore:

\begin{enumerate}
\item define the two quantities precisely;
\item prove they are equal for every $N$;
\item then push the exact computations from $N\le 25$ to $N\le 45$ by an exact meet-in-the-middle search.
\end{enumerate}

\subsection*{5) WORK (ROUND-2)}
Let
\[
\delta_{\mathrm{old}}(N)
:=
\min\left\{
\left|1-\sum_{n\in A}\frac1n\right|:
A\subseteq\{1,\ldots,N\},
\sum_{n\in A}\frac1n\neq 1,
\text{ and no strict }S\subsetneq A\text{ has }\sum_{n\in S}\frac1n=1
\right\}.
\]
Let also
\[
\delta_{\mathrm{simp}}(N)
:=
\min\left(\left\{
\left|1-\sum_{n\in A}\frac1n\right|:
A\subseteq\{1,\ldots,N\}
\right\}\cap (0,\infty)\right).
\]
The modern webpage uses $\delta_{\mathrm{simp}}(N)$, while the original Erd\H{o}s--Graham book used the old formulation. (Whether one allows $S=A$ in the old condition is immaterial, because the minimum is taken over \emph{nonzero} distances.)

\medskip
Clearly
\[
\delta_{\mathrm{simp}}(N)\le \delta_{\mathrm{old}}(N),
\]
since the old admissible sets form a subset of all sets.

\medskip
\textbf{Lemma 5.1.}
If $A\subseteq\{1,\ldots,N\}$ contains a strict subsubset $S\subsetneq A$ with
\[
\sum_{n\in S}\frac1n=1,
\]
then every nonzero distance satisfies
\[
\left|1-\sum_{n\in A}\frac1n\right|\ge \frac1N.
\]

\emph{Proof.}
Since $S\subsetneq A$, the difference $A\setminus S$ is nonempty, and every reciprocal in it is at least $1/N$. Also
\[
\sum_{n\in A}\frac1n-1=
\sum_{n\in A\setminus S}\frac1n.
\]
Hence, if the distance is nonzero,
\[
\left|1-\sum_{n\in A}\frac1n\right|
=
\sum_{n\in A\setminus S}\frac1n
\ge \frac1N.
\]
\hfill $\square$

\medskip
\textbf{Lemma 5.2.}
For every $N\ge 3$ one has
\[
\delta_{\mathrm{simp}}(N)<\frac1N.
\]

\emph{Proof.}
Use the following explicit sets.

\begin{itemize}
\item For $3\le N\le 5$, take $A=\{2,3\}$. Then
\[
\sum_{n\in A}\frac1n=\frac56,
\qquad
\left|1-\sum_{n\in A}\frac1n\right|=\frac16<\frac1N.
\]
\item For $N=6$, take $A=\{2,3,5\}$. Then
\[
\sum_{n\in A}\frac1n=\frac{31}{30},
\qquad
\left|1-\sum_{n\in A}\frac1n\right|=\frac1{30}<\frac16.
\]
\item For $N\ge 7$, take $A=\{2,5,6,7\}$. Then
\[
\sum_{n\in A}\frac1n=\frac{106}{105},
\qquad
\left|1-\sum_{n\in A}\frac1n\right|=\frac1{105}<\frac1N.
\]
\end{itemize}
This proves the claim. \hfill $\square$

\medskip
\textbf{Theorem 5.3.}
For every $N\ge 1$,
\[
\delta_{\mathrm{old}}(N)=\delta_{\mathrm{simp}}(N).
\]

\emph{Proof.}
For $N=1,2$ the equality is immediate by inspection. Assume now $N\ge 3$, and let $A$ be a set attaining $\delta_{\mathrm{simp}}(N)$; such an $A$ exists because there are only finitely many subsets of $\{1,\ldots,N\}$.
By Lemma 5.2,
\[
\left|1-\sum_{n\in A}\frac1n\right|=\delta_{\mathrm{simp}}(N)<\frac1N.
\]
By Lemma 5.1, $A$ cannot contain a strict subsubset summing to $1$. Hence $A$ is admissible for the old problem as well, so
\[
\delta_{\mathrm{old}}(N)\le \delta_{\mathrm{simp}}(N).
\]
Together with the obvious reverse inequality, this gives equality. \hfill $\square$

\medskip
\textbf{Exact computation for $26\le N\le 45$.}
By Theorem 5.3 we may compute the original quantity in the simplified model. Let
\[
L_N=[1,\ldots,N].
\]
Then every reciprocal sum can be written as
\[
\sum_{n\in A}\frac1n = \frac{1}{L_N}\sum_{n\in A}\frac{L_N}{n}.
\]
So $\delta(N)$ is exactly the minimal nonzero value of
\[
\frac{1}{L_N}\left|L_N-\sum_{n\in A}\frac{L_N}{n}\right|.
\]
I computed this by a meet-in-the-middle search on the integer weights $L_N/n$, with all arithmetic exact and no floating point. The values beyond Round 1 are:

\[
\begin{array}{c|c|l}
N & \delta(N) & \text{one witness set }A \\
\hline
26 & \dfrac{319}{608435100} & \{4,7,9,13,14,15,17,19,21,23,25,26\} \\
27\le N\le 31 & \dfrac{1}{12209400} & \{2,6,17,18,19,21,24,25,27\} \\
32 & \dfrac{7}{125215200} & \{2,11,16,17,18,20,24,25,27,31,32\} \\
33\le N\le 36 & \dfrac{1}{33599412} & \{2,7,9,17,22,23,28,31,33\} \\
37 & \dfrac{353}{79242587424} & \{2,13,16,17,19,21,23,24,32,33,36,37\} \\
38\le N\le 40 & \dfrac{1891}{1188638811360} & \{4,7,13,16,19,20,22,23,26,27,32,33,34,35,36,37,38\} \\
41\le N\le 43 & \dfrac{17}{29007778800} & \{3,5,13,16,19,22,25,26,28,30,31,40,41\} \\
44\le N\le 45 & \dfrac{167}{348985892640} & \{2,19,22,23,24,28,30,31,32,33,34,35,40,41,42,44\}
\end{array}
\]

Each displayed witness has distance $<1/N$, so by Theorem 5.3 it is automatically admissible for the original Erd\H{o}s--Graham formulation.

\subsection*{6) ADVERSARIAL VERIFICATION}
\begin{itemize}
\item The equivalence proof depends on the word ``strict'' in Lemma 5.1. If one also forbids $S=A$, nothing changes because the minimized distance is explicitly required to be nonzero.
\item The threshold $1/N$ is sharp for this argument: once a strict $1$-summing subset is present, every additional reciprocal contributes at least $1/N$.
\item The exact computations are rigorous because, after multiplying by $L_N$, the problem becomes an integer subset-sum problem. The meet-in-the-middle search inspects the exact integer distance to the target $L_N$.
\item I have not proved any asymptotic beyond the already known upper and lower bounds; the advance here is the exact equivalence of formulations and the extension of exact data from $N\le 25$ to $N\le 45$.
\end{itemize}

\subsection*{7) FINAL}
\textbf{UNRESOLVED (BUT STRICTLY ADVANCED).}

The strict advances over Round 1 are:

\begin{itemize}
\item the original Erd\H{o}s--Graham formulation and the modern simplified formulation are now proved to be \emph{exactly equal for every $N$};
\item this is done without any brute-force small-$N$ equivalence check;
\item exact values of $\delta(N)$ have been extended from $N\le 25$ to $N\le 45$.
\end{itemize}

The main asymptotic question
\[
\delta(N)=e^{-(c+o(1))N}
\]
remains open.

\subsection*{8) COMPLETION ESTIMATE}
COMPLETION: 68\%

\subsection*{9) REFERENCES}
\begin{enumerate}
\item P. Erd\H{o}s and R. Graham, \emph{Old and New Problems and Results in Combinatorial Number Theory}, problem source for \#311.
\item T. F. Bloom, \emph{Erd\H{o}s Problem \#311}, Erd\H{o}s Problems website, accessed 2026-03-07.
\item V. Kovac, comment on the Erd\H{o}s Problem \#311 discussion thread (17 Aug 2025), noting the equivalence of the two formulations.
\item Q. Tang, comments / notes linked from the Erd\H{o}s Problem \#311 discussion thread (Jan 2026) for the current best published-in-thread upper bound context.
\end{enumerate}

\subsection*{10) OUTPUT FORMAT}
This file is already in drop-in LaTeX format with no preamble.


